% ============================================================================
% 社会科学小故事 · 课堂教学材料
% ============================================================================
% 本文件为课堂讲义主体，与 syllabus.tex（课程大纲）配合使用。
% syllabus.tex 定义整体规划和课次安排，
% classroom_materials.tex 包含具体课堂内容。
% ============================================================================

\documentclass[a4paper,12pt]{ctexart}
\usepackage{xcolor}
\usepackage[top=2cm, bottom=2cm, left=2cm, right=2cm]{geometry}
\usepackage{hyperref}
\hypersetup{colorlinks=true,linkcolor=blue!70!black}
\linespread{1.5}

\title{\textcolor{blue!70!black}{\Huge\textbf{社会科学小故事}}}
\author{}
\date{}

\begin{document}
\maketitle
\thispagestyle{empty}

\newpage

\section{铅笔的故事——看不见的手}

你手里拿着一支最普通的铅笔。黄色的笔杆，顶着一块小小的橡皮，
笔杆上印着几个看不清的字母。

你有没有想过——这支铅笔是怎么来到你手上的？

\subsection{一支铅笔的"身世"}

让我们从最外面开始拆。

那圈金色的金属箍——是用铝做的。铝是从铝土矿里提炼出来的。
铝土矿在地球深处的矿井里，需要工人把它挖出来，用大卡车运到
冶炼厂，经过高温电解变成铝锭，再轧成薄片，最后冲压成一个个
小金属圈。

橡皮头——橡胶做的。橡胶树长在热带，比如印度尼西亚和泰国。
工人在橡胶树上割一个口子，白色的胶乳流出来，收集、凝固、
硫化、压制成型，切成小块——然后船运到中国的铅笔厂。

笔杆——雪松木做的。雪松长在美国俄勒冈州的森林里，需要
几十年的生长期。伐木工砍倒树，锯成木板，烘干、切割成铅笔
形状的木条。每根木条上刻出凹槽——放笔芯用的。

笔芯——石墨和黏土混合烧制的。石墨矿可能来自斯里兰卡或
巴西。矿工在地下几百米的地方把石墨矿石采出来，粉碎、提纯，
和黏土按比例混合，压成细条，在高温窑炉里烧硬。

你手里的铅笔——是全世界成千上万的人共同完成的。林场工人、
矿工、炼铝工人、橡胶种植园主、卡车司机、码头工人、船长、
工厂工人……他们来自不同的国家、说不同的语言、信仰不同的
宗教，很多人一辈子都没见过面。

\begin{center}
\fcolorbox{blue!50!black}{blue!5}{%
  \begin{minipage}{0.85\textwidth}
  \textbf{\textcolor{blue!60!black}{想想看:}} \\
  有一个人下令"做铅笔"了吗？没有。
  有一个人从头到尾把它做出来了吗？也没有。
  那这些互不认识的人为什么都在为这支铅笔出力？
  \end{minipage}
}
\end{center}

\subsection{没有人指挥的乐队}

这个问题的答案，就是经济学最核心的概念——\textbf{市场}。

市场不是一个地方，而是一种"信号机制"。价格就是信号。

当铅笔的价格足够高，高到生产铅笔能赚钱——就有人愿意去
砍木头、挖石墨、炼铝、做橡皮。他们不是为了帮你做铅笔，
他们是为了赚钱养家。但结果就是：一支铅笔被制造出来了，
你花几毛钱就能买下来。

亚当·斯密在1776年把这个现象叫做"看不见的手"——
每个人的行动都是为了自己的利益，但市场这只看不见的手
把这些分散的努力组织起来，服务于全社会的需求。

\begin{center}
\fcolorbox{green!50!black}{green!3}{%
  \begin{minipage}{0.85\textwidth}
  \textbf{\textcolor{green!50!black}{小实验:}} \\
  下一次你去超市，选一件最便宜的东西——比如说一包方便面。
  猜一猜：有多少人参与了它的生产？小麦种植者、面粉加工
  工人、棕榈油种植园工人、盐矿工人、包装印刷工人、卡车
  司机、货架理货员……你会发现，"没人管"比"有人管"
  更有效率地组织了全世界的人为你工作。
  \end{minipage}
}
\end{center}

\subsection{为什么这对你重要}

理解"看不见的手"，你就明白了三件非常重要的事：

\textbf{第一，交易让双方都受益。} 你给商店两块钱，商店给你
一瓶水。你觉得水值两块钱（甚至更多），商店觉得两块钱比那
瓶水值钱——两个人都得到了自己更想要的东西。交易不是"你赢
我输"，而是"双赢"。

\textbf{第二，分工提高效率。} 如果一个人要从砍树开始，自己
制造一支铅笔——他这辈子可能都做不出来。但因为每个人只做
自己最擅长的一件事，所有人加在一起就能生产出数以亿计的铅笔。

\textbf{第三，价格是信息。} 如果铅笔突然涨价了——说明要么
做铅笔的材料变贵了，要么做铅笔的人变少了。价格涨了，就会
有更多的人开始生产铅笔；价格跌了，生产的人就会减少。价格
就像一只温度计，告诉我们什么东西稀缺、什么东西充足。

\begin{center}
\fcolorbox{purple!50!black}{purple!3}{%
  \begin{minipage}{0.85\textwidth}
  \textbf{\textcolor{purple!60!black}{延伸思考:}} \\
  如果"看不见的手"这么厉害，是不是政府什么都不用管？
  如果市场能把一切都安排好，为什么还需要法律？还需要
  警察？还需要学校？——这就是后面的故事要告诉你的。
  \end{minipage}
}
\end{center}

\subsection*{练一练}

\begin{enumerate}
  \item 在你家里找一件最普通的物品（比如一本书、一把椅子、
  一双筷子）。试着列出至少三种参与生产的人。
  \item 为什么同样的铅笔，在学校门口的文具店比在大型超市
  贵5毛钱？用"供需"的方式解释一下。
  \item 你觉得"看不见的手"有没有做不到的事情？如果有，
  请举例。
\end{enumerate}

\end{document}
